\section{Review of the theory of Cartan geometries}\label{section:review}

\subsection{Definition}

In this section, we provide a complete definition of Cartan geometry and of parabolic geometry.
For introductions to Cartan geometries, see \cite[Chapter~1]{Cap/Slovak:2009}, \cite{mckay2023introduction,Sharpe:2002}.
Suppose that $G$ is a~complex Lie group acting holomorphically and transitively on a complex manifold $X$;
the pair~$(X, G)$ is called a \emph{complex homogeneous space}.
Let $H \subset G$ be the stabilizer of a point~$x_0 \in X$.
A holomorphic $\pr{X, G}$-geometry, also known as a holomorphic \emph{Cartan geometry} modelled on~\pr{X, G}, on a manifold $M$ is a
choice of holomorphic principal $G$-bundle $E_G \to M$ with a holomorphic connection
$\omega$ and holomorphic $H$-subbundle $E \subset E_G$ so that the tangent spaces of $E$ are complementary to the horizontal
spaces of the connection $\omega$ \cite{Sharpe:2002}.
The \emph{Cartan connection} is the restriction to $E$ of the connection $1$-form $\omega$ on $E_G$.
A Cartan geometry is \emph{effective} if the action of $G$ on $X$ is so.

The principal $H$-bundle $G \to X$ defined by $g \longmapsto gx_0$ is a Cartan geometry, called the \emph{model
geometry}, with the left invariant Maurer--Cartan $1$-form $g^{-1} {\rm d}g$ on $G$ as Cartan connection.

If $X' \subset X$ is a connected component and $G' \subset G$ is the subgroup preserving $X'$, then any \pr{X, G}-geometry
is precisely an \pr{X', G'}-geometry, so we can assume, without loss of generality, that the model $X$ of any Cartan geometry is connected.
A \emph{flag variety} is a complex homogeneous space $(X, G)$ for which $X$ is a connected and simply connected projective variety
and $G$ is a~complex semisimple Lie group; they are precisely the rational homogeneous projective varieties.
A \emph{parabolic geometry} is a Cartan geometry whose model is a flag variety \cite{Cap/Slovak:2009}.

\subsection{Dropping}
In this section, we will see that each Cartan geometry gives rise, by a trivial construction, to another one, on a higher-dimensional manifold, which we call the \emph{lift} \cite[Section 9]{mckay2023introduction}.
If a Cartan geometry arises from the lift of another, we say it \emph{drops} to that other Cartan geometry.

Suppose that $H \subset H' \subset G$ are closed complex subgroups of a complex Lie group $G$, with associated homogeneous
spaces $X := G/H$ and $X' := G/H'$, and that $(E, \omega) \to M'$ is an \pr{X', G}-geometry.
If we define $M := E/H$, then $E \to M$ is an \pr{X, G}-geometry with the same Cartan connection $\omega$;
we call it the \emph{lift} of $(E, \omega) \to M'$ \cite[Section~1.5.13]{Cap/Slovak:2009}.
(The lift is traditionally called the \emph{correspondence space} of $(E, \omega) \to M'$ \cite[Section~1.5.13]{Cap/Slovak:2009}, but this is an awkward fit to our applications of the idea, as we will see).
Note that $M \to M'$ is a holomorphic fiber bundle
with fiber $H'/H$, the same fiber as the $G$-equivariant holomorphic map $X \to X'$.
Conversely, a given $\pr{X, G}$-geometry \emph{drops} to a given $\pr{X',G}$-geometry if it is isomorphic to the $\pr{X, G}$-lift
of that $\pr{X', G}$-geometry.
(The notion of \emph{dropping} is fundamental to all of our work in this paper, more so than lifting. It would be awkward to use the expression \emph{is isomorphic to the correspondence space of} instead of \emph{drops to}.
Therefore, we adopt the terminology of \emph{lifting} and \emph{dropping} henceforth.)

\begin{Example}
If $X$ is a connected homogeneous $G$-space and $*$ is a point with trivial $G$-action, then an $(X, G)$-geometry on a
connected manifold drops to a $(*, G)$-geometry just when the geometry is isomorphic to its model.
\end{Example}

More generally, if a geometry on some complex manifold $M$ drops to some complex manifold~$M'$, then we can recover the complex manifold $M$ and the original geometry on $M$ directly from the geometry on $M'$; see \cite{Biswas.McKay:2016} for details and \cite{McKay2011b} for examples.
A \emph{minimal geometry} is a~holomorphic Cartan geometry which does not drop to a holomorphic Cartan geometry on any lower-dimensional manifold.

\subsection{From modules to vector bundles}
Any Cartan geometry is modelled on a homogeneous space.
In this section, we explain that various vector bundles naturally arise in any Cartan geometry, and that these vector bundles are modelled on homogeneous vector
bundles \cite[Section 10]{mckay2023introduction} on the homogeneous space.
Take a~Cartan geometry $E \to M$ modelled on a complex homogeneous space $(X, G)$ with~${X = G/H}$.
The principal $H$-bundle $E \to M$ and the Cartan connection determine the Cartan geometry,
as $E_G := \prodquot{E}{G}{H}$ and the Cartan connection extends uniquely to $E_G$ to become a holomorphic connection $1$-form.
To any holomorphic $H$-module $V$ we associate the holomorphic vector bundle~$\vb{V} := \prodquot{E}{V}{H} \to M$
whose global sections are the holomorphic $H$-equivariant maps~${E \to V}$.
We use the same symbol $\vb{V}$ for the associated vector bundle on the model $X$ as well.

\begin{Example} If $V := \LieG/\LieH$ equipped with the adjoint action of $H$, then $\vb{V} = \vbTM = TM$
\cite[Proposition 1]{mckay2023introduction}, \cite[Theorem 3.15]{Sharpe:1997}.
\end{Example}

\subsection{Development}\label{subsection:Development}
Suppose that $M$ is a complex manifold with a holomorphic Cartan geometry.
In this section, we consider all holomorphic maps to $M$, from any complex manifold, on which the pullback curvature of the Cartan geometry vanishes.
We call these maps \emph{pancakes}.
Such a map, in some sense, locally resembles a map to the model of the Cartan geometry, its \emph{development}.
All material in this section is explained in more detail and with complete proofs
in \cite[Section 7]{mckay2023introduction}.

Take a complex homogeneous space $(X, G)$ with marked point $x_0 \in X$.
Take a holomorphic~$(X, G)$-Cartan geometry $H \longrightarrow
E \to M$ with Cartan connection $\omega$.
A \emph{pancake} of the Cartan geometry is a holomorphic map $f \colon Z \longrightarrow
 M$ from a connected complex space so that the holomorphic connection $\omega$ on $E_G := \amal{E}{H}{G}$ pulls back to a flat
connection on $f^* E_G \to Z$. Pick a point $z_0 \in Z$.
\begin{Example}
Every holomorphic map $Z\to X$ of any complex manifold $Z$ to a homogeneous space $(X,G)$ is a pancake for the model Cartan geometry on $X$.
So maps $Z\to M$ which are not pancakes, roughly speaking, ``look different'' locally than maps $Z\to X$, hence the pancake condition.
\end{Example}
There is a unique unramified Galois covering
\[
p_Z \colon \ \big(\widehat{Z}, \widehat{z}_0\big) \to (Z, z_0)
\]
so that the flat connection has trivial holonomy precisely on covering spaces of $\widehat{Z}$.
Each choice of point \smash{$e_0 \in \pr{f \circ p_Z} ^* E_{\widehat{z}_0} = E_{z_0}$} gives an $H$-equivariant morphism $\widehat{f} \colon
 \pr{f \circ p_Z}^*E \to G$ of complex spaces, uniquely determined by
$
\widehat{f}\of{e_0} = 1
$
and $\widehat{f}^* g^{-1} {\rm d}g = \omega$ \cite[Section 2.3]{Biswas.McKay:2016}. In fact~$e_0$ produces an
isomorphism of $\pr{f \circ p_Z}^*E_G \to \widehat{Z}\times G$ of principal $G$-bundles, and
$\widehat{f}$ is the composition of maps
\[
\pr{f \circ p_Z}^*E \hookrightarrow \pr{f \circ p_Z}^*E_G = \widehat{Z}\times G \longrightarrow G.
\]
By $H$-equivariance, $\widehat{f}$ descends to a morphism
\begin{equation}\label{ed}
\widehat{f} \colon \ \big(\widehat{Z}, \widehat{z}_0\big)
\to (X, x_0) = (G/H, eH)
\end{equation}
of complex spaces; it is known as the \emph{developing map}.

Let
\[
\pi := \text{Gal}(p_Z) = \pi_1(Z)/\pi_1\bigl(\widehat Z\bigr)
\]
be the Galois group. The flat connection has holonomy morphism
$
h \colon \pi \to G
$
uniquely determined by $\widehat{f}\circ\gamma = h(\gamma)\widehat{f}$ for every $\gamma \in \pi$, where $\widehat{f}$ is the map in \eqref{ed}.
We have an exact sequence of group morphisms
\[
\begin{tikzcd}
1\ar[r]& \pi_1\bigl(\widehat{Z}, \widehat{z}_0\bigr) \ar[r]&\fundamentalGroup{Z, z_0}\arrow{r}{h}&G.
\end{tikzcd}
\]
The developing pair $\widehat{f}, h$ is \emph{associated} to the pancake $f \colon Z \to M$.
The pancake \emph{develops to the model} if $Z = \widehat{Z}$, i.e., if $\pi = 1$.

The group $\pi$ acts on $\widehat{f}^* G$ on the left by the holonomy morphism,
where $G$ is considered as a principal $H$-bundle on $X$ using the quotient map
$G \longrightarrow G/H$, and $\widehat{f}$ is the map in \eqref{ed}; quotienting: $f^*E \cong \big(\widehat{f}^*E\big)/\pi$, and $g^{-1} {\rm d}g$ descends to define a connection on $f^*E_G$, which actually coincides with $f^*\omega$. The developing pair
\[
\big(\widehat{f}, h\big), \qquad
\widehat{f} \colon \big(\widehat{Z}, \widehat{z}_0\big) \to (X, x_0),\qquad
h \colon\ \pi \to G
\]
together with choice of points $Z \ni z_0 \longmapsto m_0 \in M$ and the Cartan connection on $M$ determines uniquely
the original map $Z \to M$.

The developing pair of a pancake is uniquely determined by choice of point $e_0 \in \! \pr{f \circ p_Z} ^* E_{\widehat{z}_0}
 \!= E_{m_0}$. If we replace $e_0$ by another point $e_0 h_0$ for some $h_0 \in H$, then
\smash{$\big(\widehat{f}, h\big)$} gets replaced by \smash{$\big(h_0^{-1} \widehat{f}, h_0^{-1} h h_0\big)$}.
If we replace $z_0$ by another point of $Z$, similarly \smash{$\big(\widehat{f}, h\big)$} gets replaced by~\smash{$\big(g^{-1} \widehat{f}, g^{-1} h g\big)$} for some $g \in G$.
Every pancake has a unique developing pair, up to this right $G$-action \cite[Section~2.3]{Biswas.McKay:2016}.

The following pullback bundles have canonical isomorphisms of bundles on $\hat{Z}$ \cite[Section~2.3]{Biswas.McKay:2016}:
\begin{gather*}
\begin{split}
&f^* E \cong \big(\widehat{f}^* G\big)/\pi \text{ as holomorphic principal $H$-bundles}, \\
&f^* TM \cong \big(\widehat{f}^* TX\big)/\pi \text{ as holomorphic vector bundles}, \\
&f^* \pr{\amal{E}{H}{\LieG}} \cong \widehat{f}^* \pr{\amal{G}{H}{\LieG}}/\pi \text{ as holomorphic vector bundles with connection}.
\end{split}
\end{gather*}
If $Z$ develops to the model, these obviously become isomorphisms of bundles on $Z$:
\begin{gather*}
\begin{split}
&f^* E \cong \widehat{f}^* G \text{ as holomorphic principal $H$-bundles}, \\
&f^* TM \cong \widehat{f}^* TX \text{ as holomorphic vector bundles}, \\
&f^* \pr{\amal{E}{H}{\LieG}} \cong \widehat{f}^* \pr{\amal{G}{H}{\LieG}} \text{ as holomorphic vector bundles with connection}.
\end{split}
\end{gather*}
We will make crucial use of these isomorphisms below.

A submanifold $Z\subseteq M$ \emph{develops to the model} if $Z$ is a pancake and the inclusion map~${Z\!\to\! M}$ develops to the model.

For more on the story of development of curves to the model from flat Cartan geometries, see \cite[p.~3]{Goldman:2010}. For development to the model of curves in non-flat Cartan geometries, see
\cite[p.~108, 1.5.17]{Cap/Slovak:2009}, \cite[p.~54, 41X]{Cartan1935}, \cite{Ehresmann:1951}, \cite[p.~368]{Sharpe:1997}. For developments of higher-dimensional submanifolds, still only to the model, but only
in flat Cartan geometries (which suffices for our applications in this paper), see
\cite[Section~4]{KobayashiOchiai:1980}. The concept of development of curves was first used to develop
curves in one geometry to curves in another, with neither being assumed to be a~homogeneous model
\cite{Rinow1932}, but this was only carried out in Riemannian geometry to compare any two Riemannian
geometries on surfaces. For development of curves between arbitrary Cartan geometries with the same model,
the only reference is \cite[Section~17]{mckay2023introduction}. For development of higher-dimensional
submanifolds, in Cartan geometries which might not be flat, still only developing to the model, the only
references are \cite[Section~2.3]{Biswas.McKay:2016}, \cite[p.~43]{mckay2023introduction}, the second of
which gives complete proofs of the results we use here.

\section{Review of algebraic geometry}\label{section:algebraic.geometry}
In this section, we recall some theorems of complex algebraic geometry.
We were not able to find all of this information in a single source, so we need to give a summary of the results that we will use.
\subsection{Line bundles}
Suppose that $M$ is a compact irreducible reduced complex space.
The \emph{Iitaka dimension} of a~holomorphic line bundle $L$ on $M$ is the maximal dimension of the images of the rational maps
\[
\rationalMap{M }{ \Proj{}\cohomology{0}{M, kL}^*}
\]
taking a point $m \in M$ to the hyperplane in $\cohomology{0}{M, kL}$ consisting of sections vanishing at $m$, for $k = 1,
2, \dots$; see \cite[p.~50]{Ueno:1975}, \cite[Definition~2.1.3]{Lazarsfeld:2004}.
If $0 = \cohomology{0}{M, kL}$ for all $k > 0$, the Iitaka dimension is defined to be $-\infty$ \cite[p.~50]{Ueno:1975}, \cite[Definition~2.1.3]{Lazarsfeld:2004}.
The line bundle $L \to M$ is \emph{big} if the Iitaka dimension of $L$ is the dimension of $M$ \cite[Definition~2.2.1]{Lazarsfeld:2004}.
A nef line bundle~${L \to M}$ is big just when $c_1(L)^n > 0$, where $n := \dimC{M}$; see \cite[Theorem~2.2.16]{Lazarsfeld:2004}.
The \emph{numerical dimension} of $L \to M$ is the largest integer $k$ for which
$0 \ne c_1(L)^k \in \cohomology{2k}{M, \R{}}$, so a nef line bundle is big precisely when its numerical dimension equals the dimension of the complex manifold $M$ \cite[Remark~2.3.17]{Lazarsfeld:2004}.
The \emph{Kodaira dimension} $\kappa_M$ of $M$ is the Iitaka dimension of the canonical bundle of $M$ see \cite[p.~65]{Ueno:1975}.
The Kodaira dimension is at most the numerical dimension; a holomorphic line bundle is \emph{abundant} if its Kodaira and numerical dimensions are equal \cite[Remark~2.3.17]{Lazarsfeld:2004}.
The space $M$ is of \emph{general type} if its canonical bundle is big \cite[Example 2.2.2]{Lazarsfeld:2004}.
Any compact complex manifold with $c_1 < 0$ has general type, as does any modification of it.
\subsection{The Iitaka fibration}
An \emph{Iitaka fibration} of $M$ is a dominant rational map $\rationalMap{M}{\Ii{M}}$, whose generic fiber is smooth and irreducible, so that subvarieties of $M$ on which $\cb{M}$ has zero Iitaka dimension lie in the fibers and, on the very general fiber, $\cb{M}$ has zero Iitaka dimension \cite[Theorem~2.1.33]{Lazarsfeld:2004}.
Any projective variety of nonnegative Kodaira dimension has a unique Iitaka fibration up to birational isomorphism
\cite[Theorem~10.3]{Iitaka:1982} or \cite[Theorem~5.10]{Ueno:1975},
and its fibers are connected \cite[p.~124]{Lazarsfeld:2004}.
If some power of the canonical bundle is spanned by global sections, then there is a unique holomorphic Iitaka fibration, and its codomain is the image of the $\cb[p]{M}$-maps
\[
M \to \Proj{}\cohomology{0}{M,\cb[p]{M}}^*
\]
for all but finitely many integers $p > 0$ \cite[p.~133]{Lazarsfeld:2004}.
\subsection{Moishezon manifolds}
In this section we explain that our theorems actually hold for a larger class of complex manifolds.
All varieties in this paper are assumed complex.
A \emph{Moishezon manifold} is a connected compact complex manifold bimeromorphic to a smooth projective variety \cite[p.~26]{Ueno:1975}.
If a Moishezon manifold $M$ admits a holomorphic Cartan geometry then $M$ is a smooth projective variety; \cite[Corollary~2]{Biswas.McKay:2016}.
If the holomorphic Cartan geometry on $M$ drops to a complex manifold $M'$ then $M'$ is also a smooth complex projective variety; \cite[Corollary~2]{Biswas.McKay:2016}.
Hence the classification of holomorphic Cartan geometries on Moishezon manifolds follows immediately from the classification on smooth projective varieties, which itself follows immediately from the classification of minimal geometries on smooth projective varieties.
Consequently, Theorem~\ref{theorem:main} holds true with \emph{smooth projective variety} replaced by \emph{Moishezon manifold}.
On the other hand, there are compact K\"ahler manifolds which are not projective and which admit holomorphic flat Cartan geometries.
For example, complex tori admit flat holomorphic projective connections.

Throughout this paper, we try to prove results in the greater generality of compact K\"ahler manifolds, but we are not always able to, since many of our results use theorems of algebraic geometry which are not understood for compact K\"ahler manifolds.
We expect the results of this paper hold with little modification for compact K\"ahler manifolds.

\subsection{Stein manifolds}
A \emph{Stein manifold} is a complex manifold for which all cohomology in all positive degrees of all coherent sheaves vanishes \cite[p.~100]{Grauert/Remmert:2004}.
A complex manifold is Stein just when it admits a~holomorphic proper embedding into complex Euclidean space \cite[Theorem~5.3.9]{Hormander:1990}.
We will need to make use the existence of such an embedding for some complex Lie groups below, so we will need to know which complex Lie groups are Stein.

\subsection{Rational curves}\label{subsection:rational.curves}
A \emph{rational curve} in a complex manifold $M$ is a nonconstant holomorphic map $\Proj{1}\longrightarrow M$ from the Riemann sphere.
A complex manifold is \emph{uniruled} if a rational curve passes through its generic point \cite[p.~181]{Kollar:1996}.
A line bundle $L \to M$ on a projective variety is \emph{nef} if $\int_C c_1(L) \ge 0$ for all complete complex algebraic curves $C$ in $M$ \cite[Definition~1.4.1]{Lazarsfeld:2004}.
A projective variety is \emph{minimal} if it has nef canonical bundle \cite[p. 2]{Cascini:2013}.
By Mori's cone theorem \cite{Cascini:2013,Mori:1982}, every smooth projective variety with no rational curves is minimal.
The \emph{abundance conjecture} claims that, on any minimal projective variety, some power of the canonical bundle is spanned by its global sections \cite{Siu:2009}.

\subsection{Rational connectivity}
A rational curve in a complex manifold is \emph{tame} if it lies in a compact irreducible component of the Douady space of deformations \cite{Douady:1966}.
In any compact complex manifold dominated by a~compact K\"ahler manifold, every rational curve is tame \cite{MR715653,MR973802}; in particular, every rational curve in any smooth projective variety is tame.
A compact complex manifold is \emph{tamely rationally connected} if two general points of it lie on a connected finite union of tame rational curves \cite{Kollar:1996}.
\begin{Example}
Any connected compact complex manifold with $c_1 > 0$ is tamely rationally connected \cite{Campana:1992}.
\end{Example}
A rational curve $f \colon \mathbb{P}^1\longrightarrow M$ in a complex manifold $M$ is \emph{ample} if $f^*TM$ is a sum of line bundles of positive degree.
A smooth projective variety is tamely rationally connected just when it contains an ample rational curve \cite[Definition 1.1, Theorems 1.9,~3.7 and~3.10]{Kollar:1996}.
A complex manifold is \emph{tame} if its every rational curve lies in a compact irreducible component of its Douady deformation space \cite[p.~2]{Biswas.McKay:2016}.

Every smooth projective variety is tame, as its Douady space is its Hilbert variety \cite{Douady:1966}.
\subsection{Affine groups}\label{subsection:affine.groups}
In this section, we invent a new class of complex Lie groups which behave enough like linear algebraic groups that our theorems will easily extend to them.
A complex Lie group $G$ is \emph{affine} if some finite index subgroup of $G$ has a morphism of complex Lie groups to a complex linear algebraic group, which is an isomorphism on Lie algebras.
\begin{Example}
Any complex linear algebraic group is affine, and hence any complex semisimple Lie group is affine.
\end{Example}
\begin{Example}
Every covering group of any complex linear algebraic group is affine.
For example, if a complex linear algebraic group is connected, it admits a universal covering group.
Often even a disconnected complex linear algebraic group has a universal covering group \cite[Theorem~18.2.1, Example~18.2.2]{Hilgert.Neeb:2012}.
\end{Example}
\begin{Example}
Products of affine are affine.
\end{Example}
\begin{Example}
Any complex Lie group which admits a holomorphic representation, injective on its Lie algebra, as unipotent complex
linear transformations, is affine \cite[Corollary 14.38]{Milne:2017}.
\end{Example}
\begin{Example}
Any complex Lie group with finite fundamental group and finitely many components is affine
\cite[Corollary 16.3.9]{Hilgert.Neeb:2012}.
\end{Example}

Suppose that $G$ is affine, say with finite index subgroup $G_0 \subset G$ having morphism $G_0 \longrightarrow
\overline{G}_0$ to a complex linear algebraic group, inducing an isomorphism on Lie algebras.
If $G \to P \to M$ is a holomorphic principal $G$-bundle, then on the finite unramified covering
space \smash{$\widehat{M} := P/G_0$}, we have a holomorphic principal $G_0$-bundle $G_0 \to P
\to \widehat{M}$, and an associated principal $\overline{G}_0$-bundle $\overline{P} := \amal{P}{G_0}{\overline{G}_0}$.
Every holomorphic connection on $P \to M$ pulls back to a holomorphic connection on $P
 \longrightarrow \widehat{M}$, and conversely by averaging over the deck transformations,
every holomorphic connection on $P \to \widehat{M}$ induces a holomorphic connection on $P \to M$.
We can replace $G_0$ by a finite index subgroup, hence assume that $\overline{G}_0$ is connected.
So, up to finite covering, the existence of a holomorphic connection, or of a holomorphic flat connection, is the same for the
original bundle as for the associated bundle with connected complex linear algebraic structure group.

\section{Rational curves and Cartan geometries}\label{section:rational.curves}
We recall one of our earlier results, which is essential for the results below:
\begin{Theorem}[{\cite[Theorem 2]{Biswas.McKay:2016}}]\label{theorem:drop}
On a smooth projective variety $M$ bearing a holomorphic Cartan geometry, the following are equivalent:
\begin{enumerate}\itemsep=0pt
\item[$(1)$]
Some holomorphic Cartan geometry on $M$ is minimal.
\item[$(2)$]
Every holomorphic Cartan geometry on $M$ is minimal.
\item[$(3)$]
$M$ contains no rational curves.
\item[$(4)$]
$M$ is not uniruled.
\item[$(5)$]
$M$ is not the total space of a holomorphic fiber bundle, with positive-dimensional flag manifold fibers, over a
minimal smooth projective variety.
\item[$(6)$]
$M$ is minimal, i.e., has nef canonical bundle.
\end{enumerate}
\end{Theorem}
Therefore, the classification of holomorphic Cartan geometries on smooth projective varieties follows from the classification of minimal geometries on minimal smooth projective varieties containing no rational curves.

\section{Review of Cartan geometries on projective varieties}\label{sec:Cartan.on.proj}
We set the stage by summarizing some theorems about Cartan geometries on smooth projective varieties.

\subsection{Positive Ricci}
In this section, we summarize the state of the art about holomorphic Cartan geometries on compact complex manifolds with positive Ricci curvature.
\begin{Theorem}[{\cite[Corollary 3]{Biswas.McKay:2016}}]
Any tamely rationally connected compact complex manifold $M$ admits a holomorphic Cartan geometry, say with model $(X, G)$, if and only if $M$ is biholomorphic to $X$, in which case the geometry is isomorphic to the model geometry and $(X, G)$ is a flag variety.
\end{Theorem}

\subsection{Ricci flat}
In this section, we summarize the state of the art about holomorphic Cartan geometries on compact complex manifolds with zero Ricci curvature.
\begin{Theorem}[{\cite{Biswas.McKay:2010}}]\label{thm:cpt.Kaehler.c.1.0}
A compact connected K\"ahler manifold $M$ with $c_1(M)=0$ admits a holomorphic Cartan geometry if and only if it has an unramified covering by a complex torus.
\end{Theorem}

\begin{Theorem}[{\cite{Biswas.Dumitrescu:2017}}]\label{theorem:BD.torus}
If the group $G$ of a complex homogeneous space $(X, G)$ is affine, then every holomorphic $(X, G)$-geometry on any complex
torus is translation invariant.
\end{Theorem}

\subsection{Semipositive Ricci}
In this section, we summarize the state of the art about holomorphic Cartan geometries on compact complex manifolds with semipositive Ricci curvature.
\begin{Theorem}[{\cite[Theorem 2]{McKay:2016}}]\label{theorem:semipos}
Suppose that $M$ is a connected compact K\"ahler manifold with~${c_1 \ge 0}$.
After perhaps replacing $M$ by a finite unramified covering space, every holomorphic Cartan geometry on $M$ is the lift of a Cartan geometry on a complex torus, and in particular~$M$ is a holomorphic fiber bundle, with flag variety fibers, over a complex torus.
\end{Theorem}


\section{The anticanonical splitting}\label{section:the.anticanonical.splitting}
\subsection{The anticanonical fiber bundle}
In this section, we define the anticanonical fibration of a complex homogeneous space $X$, which is a fibration invariant under automorphisms; the anticanonical fibration was previously considered by~\cite{Huckleberry:1986}, but the results below are new.
We prove that every subvariety of a complex homogeneous space on which the ambient canonical bundle is trivial lies in a fiber of the anticanonical fibration.
We will subsequently apply this to subvarieties which are developed to $X$ from some Cartan geometry.

For any normal complex space $X$ and holomorphic line bundle $L \to X$, let
\[
V := \cohomology{0}{X, L},
\]
which might be an infinite-dimensional complex vector space.
Let $X' \subset X$ be the points of $X$ where every holomorphic section of $L$ vanishes.
If $\dimC{ V} > 0$, define the $L$-map
\[
X\setminus X' \longrightarrow \mathbb{P}V^*, \qquad x \longmapsto \{s \in V\mid s(x)=0\}.
\]
If $X' \subset X$ contains no component of $X$, denote the image of the $L$-map by $\Iitaka[L]{X}$ and the $L$-map
by~$\rationalMap[\Iitaka[L]{}]{X}{\Iitaka[L]{X}}$.
For any two points $x_0, x_1 \in X\setminus X'$ in the same fiber, we can pick a~vector~${v_0 \in L_{x_0}}$, and a global
section $s$ with $s(x_0) = v_0$, and hence define $v_1 := s(x_1)$, a natural linear isomorphism $L_{x_0}
 \cong L_{x_1}$, i.e., $L$ is trivial along the fibers of $\rationalMap[\Iitaka[L]{}]{X}{\Iitaka[L]{X}}$.
For every complex homogeneous space \pr{X, G} and $G$-invariant line bundle $L \to X$, the map $\rationalMap[\Iitaka[L]{}]{X}{\Iitaka[L]{X}}$ is $G$-equivariant.
The map is everywhere defined, by $G$-equivariance, as long as $L\to X$ has at least one holomorphic section not everywhere vanishing.
The fibers of $X \to \Iitaka[L]{X}$ are the equivalence classes: $x \sim x'$ just when every holomorphic section of $L$ which vanishes at $x$ also vanishes at $x'$ and vice versa.
The sections vanishing at $x$ have common vanishing locus a subvariety $Z_x$, and each fiber $F = X_{\overline{x}_0}$ is
\[
F = \bigcap_{x \in F} Z_x,
\]
a closed subvariety acted on by the normalizer of $H \subset G$.
The fiber $F$ is a closed complex analytic subvariety, so the stabilizer of $F$ is a closed complex subgroup $\Iitaka[L]{H}
 \subset G$, and therefore
 \[\Iitaka[L]{X} = G/\Iitaka[L]{H}\]
 is a complex homogeneous space and $X \to \Iitaka[L]{X}$ is the \emph{$L$-fiber bundle}.
The line bundle $L$ is pulled back by $\rationalMap[\Iitaka[L]{}]{X}{\Iitaka[L]{X}}$ from a line bundle on $\Iitaka[L]{X}$.
In particular, we will consider $L = \acb{X}$, i.e., $X \longrightarrow \acf{X}$, the \emph{anticanonical fiber bundle}.
Note that $X$ has global nonzero sections of its anticanonical bundle, as we can wedge together vector fields from the $G$-action.
In particular, the anticanonical fiber bundle is a well defined regular morphism.
\begin{Example}
Some easy examples, which essentially exhaust the complex homogeneous spaces in low dimension \cite{McKay:2015}:
\begin{enumerate}\itemsep=0pt
\item[(1)]
Any flag variety $X = G/P$ has positive anticanonical bundle, so $\acf{X} = X$.
\item[(2)]
If $X = \Proj{1} \times \C \times C$ where $C = \C/\Lambda$ is an elliptic curve, then $\acf{X}
 = \Proj{1} \times \C$ giving the fiber bundle $C \to X \to \acf{X}$.
\item[(3)]
If $(X, G)$ is a homogeneous Hopf manifold of dimension $n$ then $\acf{X} = X/\C^{\times} = \Proj{n-1}$.
To be precise,
\[X = (\C^{n}-\{0\})/(z\sim\lambda z)\]
 for some $\lambda \in \C^{\times}$ with $|\lambda| \ne 1$ (which
can be chosen arbitrarily to construct a homogeneous Hopf manifold), and $H^0(X, -K_X)$ consists of the tensors
\[
p(z)\partial_1\wedge\dots\wedge\partial_n
\]
for any polynomial $p(z)$ homogeneous of degree $n$.
Such a tensor, by homogeneity, vanishes on a set of complex lines through the origin in $\C^{n}$.
Hence the anticanonical fibers are the quotients of these complex lines, i.e., the fibers of the Hopf fibration $X \longrightarrow \acf{X}
 = \Proj{n-1}$, and are elliptic curves $\C^{\times}/\langle \lambda\rangle$.
\item[(4)]
The total space $X = \OO[n]$ of the line bundle $\OO[n] \to \Proj{1}$ for $n \ge 1$ has $\acf{X} = X$.
\item[(5)]
There is one nontrivial holomorphic fiber bundle $\C^{\times} \to X \to \C^{\times}$, and it admits infinitely many holomorphic faithful group actions of connected complex Lie groups \cite{McKay:2015}.
Its anticanonical fibration is $\C^{\times} \to X \to \C^{\times} = \acf{X}$.
\end{enumerate}
\end{Example}

\begin{Lemma}\label{lemma:all.points}
For any complex homogeneous space \pr{X, G}, any vector field on $X$ arising from an element of the Lie algebra of $G$ vanishes
at a point $x \in X$ just when it vanishes at all points of the anticanonical fiber through $x$.
\end{Lemma}

\begin{proof}
Denote by $\LieG$ the Lie algebra of $G$.
Pick a vector field $v \in \LieG$ and points $x, y \in X$ so that~${v(x) = 0}$ and $v(y) \ne 0$.
Let $v_1 := v$.
Since $G$ acts transitively on $X$, we can find vector fields $v_2, v_3, \dots, v_n$ so that $v_1, v_2,
 \dots, v_n$ are linearly independent at $y$.
By slight perturbation, we can arrange that $v_2, v_3, \dots, v_n$ are linearly independent at $x$.
Then $v_1 \wedge \dots \wedge v_n$ vanishes at~$x$ but not $y$.
So $x$ and $y$ lie in different fibers of the anticanonical fibration.
\end{proof}

Consequently, if we take a point $x \in X$ then its anticanonical fiber $F$ has vector fields
\[
\LieG \longrightarrow \LieG/\LieG^x \subset \cohomology{0}{F, TX}, \qquad
v \longmapsto v|_F
\]
giving a canonical trivialization of $ TX|_F$ along each anticanonical fiber $F$.
A vector field from~$\LieG$ vanishes at $x_0$ just when it belongs to $\LieH$ and just when it vanishes on the anticanonical
fiber~${F = \acf{H}x_0}$.
So $\LieH$ is precisely the subalgebra of $\LieG$ acting trivially on $\acf{H}x_0 = \acf{H}/H$.
Denote by $\acf\LieH$ the Lie algebra of $\acf{H}$; so $\LieH \subset \acf\LieH$ is an ideal.

The elements of $\acf{H}$ leaving every point of $F$ fixed form a normal subgroup of $\acf{H}$, containing $H^0$ and having Lie algebra $\LieH$.
So this normal subgroup has identity component $H^0$, and hence $H^0 \subseteq \acf{H}$ is a normal subgroup.
Thus
\[
H/H^0 \subset \acf{H}/H^0
\]
is a discrete group, giving the fiber $F = \acf{H}x_0$ the form
\[
F = \pr{\acf{H}/H^0}/\pr{H/H^0}
\]
of a quotient space of $\acf{H}/H^0$ by the action of a discrete group $H/H^0$.
In particular, the vector fields $\acf\LieH$ are a canonical $G$-invariant trivialization of the tangent bundle of each anticanonical fiber.

If $f \colon X \to \C$ is a holomorphic function with $f(x_0) \ne 0$ and $s$ is a section of the anticanonical bundle not vanishing at some point $x_0 \in X$, then $fs$ is another such section, so $f\ne 0$ on the fiber of $x_0$.

\begin{Lemma}
Take a complex homogeneous space $X = G/H$.
The identity component $\big(\acf{H}\big)^0 \subset \acf{H}$ lies in the identity component $N^0$ of the
normalizer $N := N_G H^0$ of the identity component $H^0 \subset H$.
\end{Lemma}

\begin{proof}
Pick a point $x \in X$ and vector fields $v_1, v_2, \dots, v_n \in \LieG$, linearly independent at the generic point
but with $v_1(x) = 0$ while $v_2, v_3, \dots, v_n$ are linearly independent at $x$.
Let
\[
s := v_1 \wedge v_2 \wedge \dots \wedge v_n.
\]
For any element $g \in \acf{H}$, since $s(x) = 0$, we must have $s(gx) = 0$.
So for any element $w \in \acf\LieH$, $\LieDer_w s$ vanishes at $x$:
\begin{align*}
\LieDer_w s
& =
\pr{\LieDer_w v_1} \wedge v_2 \wedge \dots \wedge v_n - v_1 \wedge \pr{\LieDer_w v_2} \wedge v_2 \wedge \dots \wedge v_n + \cdots,
\end{align*}
so that at $x$:
\[
0 = \pr{\LieDer_w s}(x) = \pr{\LieDer_w v_1}(x) \wedge v_2(x) \wedge \dots \wedge v_n(x).
\]
Since $v_2, \dots, v_n$ can be generically chosen, $\LieDer_w v_1$ vanishes at $x$.
Hence $\LieDer_w \LieH \subset \LieH$, i.e., $w$ is in the Lie algebra of $N_{\LieG} \LieH$.
So $\acf\LieH \subset N_{\LieG} \LieH$.
\end{proof}

\begin{Corollary}\label{corollary:trivial.conn}
Take a complex homogeneous space $X = G/H$.
In the holomorphic $H$-mod\-ule~$\LieG/\LieH$, the subspace $\acf\LieH/\LieH$ is acted on trivially by $H^0$, so the action
of $H$ on $\acf\LieH/\LieH$ descends to an action of the group of components $H/H^0$ on $\acf\LieH/\LieH$.
\end{Corollary}

\begin{proof}
Every element of $\acf\LieH$ lies in the normalizer $N_{\LieG} \LieH$, so for any $v \in \LieH$ and $w \in
\acf\LieH$, $[v, w] \in \LieH$.
Hence $\LieH$ acts trivially on $\acf\LieH/\LieH$ in the adjoint representation on $\LieG/\LieH$.
\end{proof}

\begin{Lemma}\label{lemma:in.the.fiber}
Take a complex homogeneous space $X = G/H$.
Then every holomorphic map $f \colon Z \to X $ from a connected normal compact complex space for
which $f^*\cb{X}$ is trivial lies in a fiber of the anticanonical fibration $X \to \acf{X}$.
\end{Lemma}

\begin{proof}
Any section of $\acb{X}$ vanishing at one point of $f(Z)$ pulls back to the zero section of~$f^*(\acb{X})$.
\end{proof}

\begin{Lemma}\label{lemma:Stein.G}
Every complex Lie group with countably many components whose identity component is simply connected is a Stein manifold.
\end{Lemma}

\begin{proof}
Clearly it is sufficient to check the identity component.
By Levi decomposition (see \cite[Theorem 16.3.7]{Hilgert.Neeb:2012}), the radical is simply connected.
The center belongs to the radical, and the radical contains no complex torus, so the center is reductive, and we apply
\cite[Theorem 1]{Matsushima/Morimoto:1960}.
\end{proof}

\subsection{Developing and anticanonical fibrations}

\begin{Lemma}\label{lemma:ab.var.in.fiber}
Take a complex homogeneous space $X = G/H$.
Take a finite holomorphic map $f \colon Z\to X$ from a connected projective variety for which $\cb{Z}$ and $f^*\cb{X}$ are trivial.
Then $f(Z)$ lies in a fiber of the anticanonical fibration $X\to \acf{X}$.
If $Z$ is smooth then it is an abelian variety, and the map $f$ is an immersion, equivariant for a morphism of complex Lie groups from a complex abelian Lie group acting transitively on the abelian variety $Z$, and the image of $f$ is an abelian subvariety of $X$.
\end{Lemma}

\begin{proof}
By Lemma~\ref{lemma:in.the.fiber}, $f$ maps to a fiber.
As above, the subgroup preserving the fiber acts transitively, with discrete stabilizer, so we can replace $X$ by that fiber, and then without loss of generality $X=G/\Gamma$ for some connected complex Lie group $G$ and discrete subgroup ${\Gamma \subset G}$.
After perhaps replacing by a finite unramified covering, any smooth projective variety $Z$ with trivial canonical bundle splits into a product $Z=A \times B$ where $A$ is an abelian variety and~$B$ is a simply connected Calabi--Yau manifold \cite{Bogomolov:1974}.
So the map $Z\to X$ lifts on each subvariety~$\{a_0\} \times B$, for any $a_0 \in A$, to a map $B
\to \widetilde{G}$ to the universal covering group of the identity component of $G$.
But that universal covering group is a Stein manifold by Lemma~\ref{lemma:Stein.G}, so embeds holomorphically in Euclidean space.

Since $B$ is a compact complex manifold, every holomorphic map of $B$ to Euclidean space is constant, by the maximum principle.
So the map $B\to\widetilde{G}$ is constant.
Hence $B$ is a point, so $Z = A$ is an abelian variety.
Trivialize the tangent bundle of $A$ with a basis of holomorphic vector fields $v_a$, and trivialize the tangent bundle of $X
 = G/\Gamma$ with a basis of holomorphic vector fields $w_i$ from the Lie algebra $\LieG$ of $G$.
The map $f \colon A \to X$ then gives relations $f_* v_a = c^i_a w_i$.
These $c^i_a \colon A \to \C$ are holomorphic functions, so constants.
In particular $f \colon A \to X$ is a~holomorphic immersion.
With a change of basis, we arrange that $v_a = w_a$.
\end{proof}

In the proof above, if $(X,G)$ is an algebraic homogeneous space, it is not clear to the authors whether $\tilde{G}$ is an affine Lie group.

\subsection{The ant fibration}\label{subsec:ant}

In this section, we define a new refinement of the anticanonical fibration, which we believe has independent interest.
Take the anticanonical fibration of the anticanonical fibration, i.e., fiberwise we look at two points of a fiber as equivalent if they have the same sections of the canonical bundle of that fiber vanishing at both.
Since the tangent bundle of each fiber is trivial, this is the same as saying that two points in the same fiber are equivalent if every holomorphic function on the fiber which vanishes at one vanishes at both.
This equivalence relation gets iterated, and once we finish, we have a homogeneous holomorphic fiber bundle with perhaps finer fibers, which we call the \emph{ant fibration}.
On each fiber of the ant fibration, all holomorphic functions are constant, and there is a unique holomorphic section of the anticanonical bundle, up to constant scaling.
Locally constant functions on any fiber are holomorphic, so the fibers of the ant fibration are connected.
Throughout our work below, the reader could make use of the anticanonical fibration instead of the ant fibration, with the advantage of being more familiar to experts in homogeneous complex manifolds, but the disadvantage of providing possibly looser restrictions on developing maps.

\begin{Example}
There is one nontrivial holomorphic fiber bundle $\C^{\times} \to X \to \C^{\times}$, and it admits
infinitely many holomorphic faithful group actions of connected complex Lie groups~\cite{McKay:2015}. Its
anticanonical fibration is $\C^{\times} \to X \to \C^{\times}$, as the holomorphic sections of the
canonical bundle are pulled back from the base of the fibration, which is easy to check from the explicit
description of $X$ in \cite{McKay:2015}. But the ant fibration is trivial $X \to X$, since the
anticanonical fiber~$\C^{\times}$ has its points separated by holomorphic functions.
\end{Example}

\begin{Lemma}\label{lemma:ant.Stein}
The universal covering space of each ant fiber is a connected and simply connected complex Lie group and hence a Stein manifold.
\end{Lemma}

\begin{proof}
Write the ant fibration as $X\to \overline{X} = G/\overline{H}$ and a fiber as \pr{X_0,\GZ} where $X_0 := \overline{H}/H$.
By Lemma~\ref{lemma:all.points}, adjoint $\overline{H}$-action preserves $\LieH$, so $H^0 \subseteq \overline{H}$ is a normal subgroup.
Hence $X_0$ is a~quotient of a complex Lie group by a discrete subgroup:
\[
X_0 = \overline{H}/H = \pr{\overline{H}/H^0}/\pr{H/H^0}.
\]
So the universal covering space of each fiber $X_0$ is a complex Lie group, and simply connected, so is a Stein manifold by Lemma~\ref{lemma:Stein.G}.
\end{proof}
Note that
\begin{align*}
&\cohomology{0}{X_0,\OO}=\C,\qquad
\cohomology{0}{X_0,T}=\LieG_0=\overline{\LieH}/\LieH,\qquad
\cohomology{0}{X_0,\acb{}}=\C.
\end{align*}

\begin{Lemma}\label{lemma:ab.var.in.fiber.2}
Take a complex homogeneous space $X = G/H$.
Take a finite holomorphic map~${f \colon Z \to X}$ from a connected projective variety for which $\cb{Z}$ and $f^*\cb{X}$ are trivial.
Then~$f(Z)$ lies in a fiber of the ant fibration.
Suppose that $Z$ is smooth.
The variety $Z$ is an abelian variety.
The map $f$ is an immersion, equivariant for a morphism of complex Lie groups from a complex abelian Lie group acting transitively on the abelian variety $Z$.
The image of $f$ is an abelian subvariety of $X$.
\end{Lemma}

\begin{proof}
By Lemma~\ref{lemma:ab.var.in.fiber}, $f$ maps to a fiber of the anticanonical fibration, so we can replace $X$ by that fiber, and $G$ by the automorphism group of that fiber, and repeat.
All holomorphic functions on $X$ pull back to constants on $Z$, so $f$ maps to a fiber of the ant fibration.
\end{proof}

\subsection{The ant foliation}
Since a Cartan geometry looks infinitesimally like its model, it bears a foliation which looks infinitesimally like the ant fibration.
In this section and the following sections we prove that the ant foliation splits the tangent bundle.

Take a holomorphic Cartan geometry $E \to M$ with model \pr{X,G} on a complex manifold~$M$.
Every holomorphic $G$-invariant vector subbundle $\vb{V} \subset TX$ has an associated
holomorphic vector subbundle $\vb{V} \subset TM$ by letting $V \subset T_{x_0} X$ be the
subspace of $\vb{V}$ at one point $x_0 \in X$, say with stabilizer $H \subset G$ and then
\[
\vb{V} := E \times^H V \subset E \times^H \pr{\LieG/\LieH}.
\]
If $\vb{F}$ is a $G$-invariant foliation of $X$, then its tangent bundle $\vb{V} := T\vb{F}$
becomes a subbundle~${\vb{V} \subset TM}$, but might not be tangent to a foliation of $M$.
The \emph{ant distribution} of $TM$ is the subbundle associated to the tangent bundle of the ant
foliation of $X$.

\begin{Lemma}
The ant distribution of any holomorphic Cartan geometry bears a flat holomorphic connection.
\end{Lemma}

\begin{proof}
Because $H^0$ acts trivially on $\acf\LieH/\LieH$, the associated vector bundle is
holomorphically trivial on $E/H^0$, a covering space.
Apply induction, taking anticanonical fibration of the anticanonical fibration.
\end{proof}

\begin{Lemma}
Take a smooth projective variety $M$ with a minimal geometry.
The ant distribution of the geometry is the tangent bundle of a foliation, the \emph{ant foliation} of the Cartan geometry, and admits a smooth holomorphic complementary foliation splitting $TM$.
\end{Lemma}

\begin{proof}
Since the ant distribution has a holomorphic connection, its Chern classes vanish.
Since~$M$ is not uniruled, every holomorphic distribution with trivial first Chern class is a smooth holomorphic foliation and admits a smooth holomorphic complementary foliation splitting $TM$ \cite[Theorem 5.2]{Loray/Pereira/Touzet:2011}.
\end{proof}

\section{Splitting}\label{section:splitting}

\begin{Lemma}\label{lemma:symmetry.splitting}
Take a smooth projective variety $M$ with a minimal geometry $E \to M$ modelled on a complex homogeneous space \pr{X,G}, $X = G/H$.
Let $V \subset \LieG/\LieH$ be the set of all vectors invariant under the $H$-action and let $\vb{V}:= E \times^H V \subset E \times^H \pr{\LieG/\LieH} = TM$.
Then $M$ splits~${M = M_0 \times A}$ where $A$ is an abelian variety and $M_0$ has no nonzero holomorphic vector fields: $0=\cohomology{0}{M_0,TM_0}$.
The distribution $\vb{V} \subset TM$ is the tangent bundle of a foliation on $M$.
Every leaf of that foliation lies inside an abelian variety $\{m_0\} \times A \subset M$.
\end{Lemma}

\begin{proof}
The distribution is $\vb{V} = E \times^H V \cong \pr{E/H} \times V$ trivial.
So each element $v \in V$ determines a nowhere vanishing holomorphic tangent vector field on $M$.
Since $M$ is minimal, i.e., contains no rational curves, $M$ is not uniruled (i.e., does not have a rational curve through the generic point).
A theorem of Liebmann \cite[Theorem 1.1]{Pereira/Touzet:2012} says that any smooth projective variety $M$ which is not uniruled splits as a product $M_0\times A$ where $A$ is an abelian variety, so that $M_0$ has zero as its only holomorphic vector field.

Since $A$ is a complex torus, we can write $A=W/\Lambda$ where $W$ is complex Euclidean space and~$\Lambda\subseteq W$ is a lattice.
Note that $TA$ is trivial: $TA=A\times W$.
Every tangent vector to $A$ extends uniquely to a nowhere vanishing holomorphic tangent vector field, a constant map to~$W$.

Since $TM=TM_0\times TA$, every holomorphic vector field $v$ on $M$ splits into a sum $v=(u,w)$ of a holomorphic vector field $u$ tangent to $M_0\times\{*\}$ and one $w$ tangent to $\{*\}\times A$.
The map
\[
p\in M\mapsto w(p)\in W
\]
maps a compact complex manifold holomorphically to Euclidean space, so is constant by the maximum principle.
Hence $v=u+w$ has $w$ constant.

For each fixed $a_0\in A$, the map
\[
p\in M_0\mapsto u(m_0,a_0)\in T_p M_0,
\]
is a holomorphic vector field, so vanishes.
Hence the holomorphic vector fields on $M$ are precisely the constant vector fields $(0,w)$ for a constant $w\in W$.
Since our distribution is spanned by global holomorphic tangent vector fields, these are of this form $(0,w)$, hence tangent to the tori~$\{*\}\times A$.
\end{proof}

\begin{Corollary}\label{cor:splitting}
Take a smooth projective variety $M$ with a minimal geometry $E \to M$ modelled on a complex
homogeneous space \pr{X,G}, $X = G/H$.
Take the ant fibration $X \to X' = G/H'$.
Suppose that $H$ has finitely many components.
Then, after perhaps replacing $M$ by a~finite unramified covering space, $M$ splits
$M = M_0 \times A$ where $A$ is an abelian variety and $M_0$ is a~smooth complex projective
variety with no nonzero holomorphic vector fields: $0 = \cohomology{0}{M_0,TM_0}$.
Every leaf of the ant foliation on $M$ lies inside an abelian variety~${\{m_0\} \times A \subset M}$.
\end{Corollary}

\begin{proof}
Replace $M$ by $E/H^0$, $H$ by $H^0$, with the same total space $E$ and the same Cartan connection, so that we can assume that $H$ is connected.
So $H$ acts trivially on $\LieH'/\LieH$.
\end{proof}

\section{Infinitesimally algebraic models}\label{section:inf.alg.models}
In this section, we define a class of complex homogeneous space which resemble complex linear algebraic homogeneous spaces.
We need a strong enough resemblance so that our theorems about Cartan geometries easily generalize to these models, even though the proofs use algebraic geometry methods.

Take a complex Lie group $G$ and some finite-dimensional holomorphic $G$-modules $V_i$.
A~\emph{reductive $($semisimple$)$ Levi quotient} of $G$ over $\{V_i\}$ is a quotient of complex Lie groups
\[
1 \to \urad{G} \to G \to \red{G} \to 1,
\]
and a complex Lie group morphism $\urad{G} \to \urad{\overline{G}}$ to a unipotent (solvable) complex
linear algebraic group, injective on Lie algebras, so that $\red{G}$ is a reductive (semisimple) complex
linear algebraic group and so that the action of $\urad{G}$ on each $V_i$ factors through a polynomial
action of $\urad{\overline{G}}$.
It follows that $\urad{G}$ preserves a filtration of $G$-modules on each $V_i$ by Lie's theorem
(see \cite[Theorem~3]{Serre:2001}), and acts trivially on the quotients $V_{i+1}/V_i$
\cite[Section 4.8]{Borel:1991}.
\begin{Example}

The group $G := \PSL{2,\C} \times \PSL{2,\Q{}}$ has no reductive or semisimple Levi quotient in the
usual sense of those terms (as defined in Vinberg \cite[p.~20]{Vinberg:1994}, for instance), nor does any finite index
subgroup \cite[Theorem 6.14]{Jacobson:1985}, but $G$ has a unique semisimple Levi quotient (also a~unique reductive Levi quotient) over $\LieG$:
\[
1 \to \PSL{2,\Q{}} \to G \to \PSL{2,\C} \to 1.
\]
\end{Example}
A complex homogeneous space $X=G/H$ is \emph{infinitesimally algebraic} if, after perhaps replacing $G$ and $H$ by finite index subgroups, $H$ has a reductive Levi quotient over $\LieG$, for which the action of $\red{H}$ on $\LieG/\LieH$ is faithful, and $G$ has a semisimple Levi quotient over $\LieG$.
\begin{Example}
Let $G$ be the set of all complex $3 \times 3$ matrices of the form
\[
\begin{pmatrix}
{\rm e}^a & 0 & b \\
0 & 1 & a \\
0 & 0 & 1
\end{pmatrix}
\]
and let $H \subset G$ be the subgroup of matrices of the form $a \in \pi {\rm i} \Z$, $b \in \Z$, i.e.,
\[
\begin{pmatrix}
(-1)^k & 0 & \ell \\
0 & 1 & \pi {\rm i} k \\
0 & 0 & 1
\end{pmatrix}
\]
for $k,\ell \in \Z$.
The quotient $X := G/H$ is the nontrivial holomorphic $\C^{\times}$-fiber bundle $\C^{\times} \to X
 \to \C^{\times}$ \cite[Section~11]{McKay:2015}.
We want to see that $(X, G)$ is infinitesimally algebraic, but not algebraic.
The adjoint action of an element
\[
g =
\begin{pmatrix}
{\rm e}^a & 0 & b \\
0 & 1 & a \\
0 & 0 & 1
\end{pmatrix}
\]
on
\[
A=
\begin{pmatrix}
\alpha & 0 & \beta \\
0 & 0 & \alpha \\
0 & 0 & 0
\end{pmatrix}
\in \LieG
\]
is
\[
\Ad_g A =
\begin{pmatrix}
\alpha & 0 & {\rm e}^a\beta-b\alpha \\
0 & 0 & \alpha \\
0 & 0 & 0
\end{pmatrix},
\]
so that in terms of vectors denoted
\[
\begin{pmatrix}
\alpha \\
\beta
\end{pmatrix},
\]
the adjoint representation factors through the solvable linear algebraic group of matrices
\[
\begin{pmatrix}
1 & 0 \\
-b & c
\end{pmatrix}
\]
with $c\ne 0$, by
\[
g \longmapsto
\begin{pmatrix}
1 & 0 \\
-b & {\rm e}^a
\end{pmatrix}.
\]
The semisimple quotient of $G$ is thus $G_{\rm ss} = 1$.
We let $\red{H} = \{\pm 1\}$, and map
\[
H \longrightarrow \red{H}, \qquad
\begin{pmatrix}
(-1)^k & 0 & \ell \\
0 & 1 & \pi {\rm i} k \\
0 & 0 & 1
\end{pmatrix}
 \longmapsto (-1)^k.
\]
The kernel $\urad{H}$ is the set of matrices of the form
\[
\begin{pmatrix}
1 & 0 & \ell \\
0 & 1 & 2\pi {\rm i} k \\
0 & 0 & 1
\end{pmatrix}
\]
for $k \in \Z$.
Hence the representation of $\urad{H}$ on $\Ad_G$ factors through
\[
\urad{H} \longrightarrow \urad{\overline{H}} := \left\{\begin{pmatrix}1&0\\-b&1\end{pmatrix}\right\}.
\]
The reader can check that the center of $G$ is the infinite discrete group
\[
\begin{pmatrix}
1 & 0 & 0 \\
0 & 1 & 2\pi {\rm i} k \\
0 & 0 & 1
\end{pmatrix},
\]
so $G$ is not isomorphic as a complex Lie group to any complex linear algebraic group.
Similarly, $H$ is discrete and infinite, so not a linear algebraic group.
As $G$ is connected, $G$ is the only finite index subgroup of~$G$.
\end{Example}

\begin{Lemma}
Suppose that $(X, G)$ is an effective connected complex homogeneous space $X=G/H$.
Suppose that $G$ has finitely many components and $H$ is isomorphic as a complex Lie group to a complex linear algebraic group.
Then $(X, G)$ is infinitesimally algebraic.
\end{Lemma}

\begin{proof}
We can assume that $G$ is connected, because $G$ has finitely many components.
We invoke the usual semisimple Levi quotient \cite[p. 20]{Vinberg:1994}.

Every complex linear algebraic group has an algebraic Levi decomposition $H = \red{H} \ltimes \urad{H}$ into a maximal reductive and its unipotent radical \cite[Theorem 4.3]{Hochschild:1981}, and so has a reductive quotient (over all $H$-modules simultaneously).
Take an $H$-invariant filtration \[0 = F_0 \subset F_1 \subset \cdots
 \subset F_N = \LieG/\LieH\] so that $\urad{H}$ acts trivially on the associated graded.
Each $F_i$ is $H$-invariant, so $\red{H}$-invariant, so $F_i$ has a $\red{H}$-invariant splitting $F_i=F_{i-1} \oplus F_{i-1}^{\perp}$.
We claim that $\red{H}$ acts faithfully on the associated graded.
Suppose that some element $g \in \red{H}$ acts trivially on the associated graded.
Note that~$g$ acts trivially on $F_1$, since $F_1$ is a subspace of the associated graded.
Suppose by induction that~$g$ acts trivially on $F_{i-1}$.
If $v \in F_i$, split as $v=u+u^{\perp}\in F_{i-1}\oplus F_{i-1}^{\perp}$.
So
\begin{align*}
gv-v
&=
g\bigl(u+u^{\perp}\bigr)-u-u^{\perp}=
gu^{\perp}-u^{\perp}.
\end{align*}
But $gv-v \in F_{i-1}$ so $gu^{\perp}-u^{\perp} \in F_{i-1}$.
But $u^{\perp} \in F_{i-1}^{\perp}$ which is $\red{H}$-invariant, so $gu^{\perp}-u^{\perp} \in F_{i-1} \cap F_{i-1}^{\perp}=0$.
So $gv = v$.
Every reductive complex linear subgroup of $H$ acts faithfully on~$\LieG/\LieH$ \cite[Lemma~6.1]{McKay2013}, so $g = I$.
\end{proof}
